\chapter{Technologie wykorzystane w projekcie}

\section{Technologie interfejsu użytkownika}

Do wykonania interfejsu użytkownika wykorzystano WPF. Technologia ta została wybrana, ponieważ dobrze nadaje się do tworzenia aplikacji desktopowych w C\# i pozwala budować czytelne okna z tabelami, formularzami oraz zakładkami.

W projekcie WPF zostało użyte do przygotowania głównego okna aplikacji, panelu statystyk, formularzy, przycisków, tabel \texttt{DataGrid}, zakładek \texttt{TabControl} oraz prostego stylu wizualnego. Interfejs został przygotowany w spokojnych kolorach niebieskich i jasnych, aby przypominał panel administracyjny recepcji.

\section{Technologie dostępu do danych}

Do obsługi danych wykorzystano Entity Framework Core. W projekcie służy on do łączenia klas C\# z tabelami w bazie danych. Dane można dodawać i pobierać za pomocą obiektów, bez ręcznego pisania wielu zapytań SQL.

Jako baza danych została użyta SQLite. Plik bazy danych ma nazwę \texttt{hotel.db}. Jest to dobre rozwiązanie dla projektu studenckiego, ponieważ nie wymaga instalacji osobnego serwera bazy danych i łatwo przenieść projekt na inny komputer.

\section{Dodatkowe biblioteki i narzędzia}

W projekcie użyto pakietów Entity Framework Core potrzebnych do pracy z SQLite i migracjami. Najważniejsze pakiety to:

\begin{itemize}
  \item \texttt{Microsoft.EntityFrameworkCore},
  \item \texttt{Microsoft.EntityFrameworkCore.Sqlite},
  \item \texttt{Microsoft.EntityFrameworkCore.Tools},
  \item \texttt{Microsoft.EntityFrameworkCore.Design},
  \item \texttt{Microsoft.Extensions.DependencyInjection}.
\end{itemize}

\section{Uzasadnienie wyboru technologii}

WPF został wybrany, ponieważ pozwala uzyskać więcej punktów za warstwę interfejsu niż Windows Forms i umożliwia wygodne tworzenie okien desktopowych. Entity Framework Core został wybrany, ponieważ jest wyżej punktowany niż ręczne ADO.NET i upraszcza pracę z bazą danych. SQLite został wybrany, ponieważ projekt można uruchomić lokalnie bez dodatkowej konfiguracji serwera.

Zastosowanie MVVM, Dependency Injection, repozytoriów i serwisów zwiększa czytelność projektu oraz ułatwia omówienie architektury podczas obrony.
